\section{Introduction}
\label{sec:intro}

Reinforcement learning from outcome rewards has become the standard recipe for training language
models to reason~\cite{shao2024deepseekmath,guo2025deepseekr1} and, increasingly, to act as agents over
many turns~\cite{feng2025gigpo,wang2025ragen,liang2026gaca}. The dominant algorithms are
\emph{critic-free}: instead of learning a value function, GRPO~\cite{shao2024deepseekmath} samples a
group of trajectories for the same prompt and standardizes each trajectory's return against the group.
The price is coarse credit. In a fifty-step household task every token of every step receives the
same advantage, whether it was the decisive choice of where to look or a routine move along the way.

Two recent methods refine this to the level of an environment step. GiGPO~\cite{feng2025gigpo}
regroups the visits that different rollouts make to the same environment state (an \emph{anchor
state}) and contrasts their discounted returns, which yields a step-level advantage without a critic.
GACA~\cite{liang2026gaca} observes that the right granularity depends on the state and mixes the
episode- and step-level advantages with a per-step weight $\lambda_t$ set by the policy's uncertainty,
the mean token negative log-likelihood (NLL) of the sampled action; it reports 96.9\% success on
ALFWorld at 1.5B parameters, against 86.7\% for GiGPO and 72.8\% for GRPO.

Both methods stop at the step. Yet a modern agent's ``action'' is not a short command: in the
reasoning format these methods are trained with, the model writes roughly a hundred tokens of
chain of thought inside \texttt{<think>} tags and then a ten-token command inside \texttt{<action>}
tags, and every one of those tokens still receives its step's single advantage. If a step was good
because of one sentence of planning, the policy gradient spends most of its credit reinforcing the
boilerplate around it. This paper asks a direct question: \emph{where inside a response is the
decision actually made, and can credit be sent there?}

We answer it in two parts. First, \textbf{measurement}. We run a counterfactual flip test on real
ALFWorld steps: cut the response after $j$ tokens, resample the remainder many times, and record how
often the final action changes. The drop of that flip rate across a segment is the share of the
decision settled inside it. This yields three findings. (i) The decision is made in the reasoning: 82\%
of it is settled before the \texttt{<action>} tag. (ii) Token entropy, the signal behind recent
token-level methods for single-turn reasoning~\cite{wang2025beyond}, does not locate it: the last tenth
of a response carries 19\% of the decision but only 3\% of the entropy, and entropy weighting moves
credit toward decisive segments by only 4\%. (iii) A cheap \emph{commitment probe} does: at each cut,
force the reasoning closed and score the probability of the eventual action under the model. This is
exactly the truncate-and-force measurement that Yu~\etal~\cite{yu2026cir} use to define Causal
Importance of Reasoning (CIR) and that Lanham~\etal~\cite{lanham2023faithfulness} call early
answering. Weighting each token by the probe's rise across its segment places $1.65\times$ more credit
on decisive segments than uniform credit at our training configuration, with the effect consistent
across all six task types, on held-out games, and on a trained policy. Along the way we find that
GACA's own per-step score is about 85\% noise from which reasoning path happened to be sampled.

Second, \textbf{training}. Yu~\etal~use CIR as a scalar bonus in the reward; every token of a response
still shares one advantage. We instead use the probe's per-segment rise to \emph{redistribute} each
step's GACA advantage over the tokens of its response, with weights that are clipped and average to
one, so a step's total credit is unchanged and only where it lands moves
(\cref{fig:method}). The probe costs a teacher-forced scoring pass and no generation. In a study on
ALFWorld with full fine-tuning of Qwen2.5-1.5B-Instruct inside the authors' own training stack
(verl-agent~\cite{feng2025gigpo,sheng2025hybridflow}), with every method evaluated on the same 128
held-out games and compared game by game, probe-weighted credit reaches the highest final success of
any arm. Our budget is about 6\% of GACA's, one seed per arm; at that budget no difference between
methods is statistically resolved, and we report it as such.

\paragraph{Contributions.}
\begin{itemize}\itemsep0pt
\item A counterfactual measurement of where decisions are made inside agent reasoning, and evidence
  that the commitment probe tracks it while entropy and NLL do not and attention only weakly
  does (\cref{sec:measure}).
\item Probe-weighted token credit: a critic-free redistribution of step-level advantage over tokens,
  compatible with GRPO, GiGPO and GACA (\cref{sec:method}).
\item A budget-matched, game-paired comparison of GRPO, GACA and five treatments on ALFWorld, with
  negative results reported (\cref{sec:results}), and a measurement that GACA's criticality score is
  dominated by within-state noise.
\end{itemize}

\begin{figure*}[t]
\centering
\resizebox{\textwidth}{!}{% Method overview: credit at three granularities. Colours follow the paper's palette:
% blue = GRPO/episode, orange = GACA/step, pink = probe/token.
\definecolor{cE}{HTML}{2A78D6}
\definecolor{cS}{HTML}{EB6834}
\definecolor{cP}{HTML}{E87BA4}
\definecolor{cInk}{HTML}{0B0B0B}
\definecolor{cMuted}{HTML}{52514E}
\definecolor{cPanel}{HTML}{F6F6F3}
\begin{tikzpicture}[
  x=1cm, y=1cm,
  tok/.style={minimum width=2.6mm, minimum height=5mm, inner sep=0, draw=white, line width=0.4pt},
  ttl/.style={font=\footnotesize\bfseries, anchor=south west, text=cInk},
  stepl/.style={font=\scriptsize, text=cMuted},
  eq/.style={font=\scriptsize, text=cInk, anchor=north west, align=left},
  panel/.style={fill=cPanel, draw=none, rounded corners=2pt},
]
% ---------------- panel backgrounds
\fill[panel] (-0.15,-1.55) rectangle (5.35,1.55);
\fill[panel] (5.75,-1.55) rectangle (11.25,1.55);
\fill[panel] (11.65,-1.55) rectangle (17.3,1.55);

% ---------------- panel 1: GRPO
\node[ttl] at (0,1.1) {(a) GRPO: one advantage per trajectory};
\foreach \s in {0,1,2} {
  \pgfmathsetmacro{\x}{\s*1.75}
  \foreach \k in {0,...,4} { \node[tok, fill=cE!60] at ({\x+0.13+\k*0.27}, 0.3) {}; }
  \node[stepl] at ({\x+0.67}, -0.15) {step $\s$};
}
\node[eq] at (0,-0.45) {$A_{t,j} = \Aep$ for every step $t$ and token $j$\\
  $\Aep$: trajectory return, standardized in its group};

% ---------------- panel 2: GiGPO / GACA
\begin{scope}[xshift=5.9cm]
\node[ttl] at (0,1.1) {(b) GiGPO / GACA: one advantage per step};
\foreach \s/\sh in {0/25,1/85,2/45} {
  \pgfmathsetmacro{\x}{\s*1.75}
  \foreach \k in {0,...,4} { \node[tok, fill=cS!\sh] at ({\x+0.13+\k*0.27}, 0.3) {}; }
  \node[stepl] at ({\x+0.67}, -0.15) {step $\s$};
}
\node[eq] at (0,-0.45) {$A_t = \lambda_t\,\Ast_t + (1-\lambda_t)\,\Aep$\\
  $\Ast_t$: contrast among visits to the same state\\
  $\lambda_t$: GACA's gate on the response's NLL};
\end{scope}

% ---------------- panel 3: probe-weighted tokens
\begin{scope}[xshift=11.8cm]
\node[ttl] at (0,1.1) {(c) Ours: the step's advantage, spread by the probe};
% one step's response: reasoning tokens then action tokens
\foreach \k/\w in {0/15,1/15,2/18,3/35,4/90,5/95,6/30,7/15,8/15,9/15,10/40,11/40,12/40,13/40} {
  \node[tok, fill=cP!\w] at ({0.13+\k*0.27}, 0.3) {};
}
\draw[cInk, line width=0.5pt] ({0.13+9.5*0.27},0.0) -- ({0.13+9.5*0.27},0.6);
\node[stepl, anchor=north] at (1.35,0.0) {reasoning};
\node[stepl, anchor=north] at (3.3,0.0) {action};
% probe curve q(j), drawn in its own band to the right of the title line
\begin{scope}[shift={(0.13,0.66)}]
  \draw[cMuted, line width=0.3pt] (0,0) -- (3.6,0);
  \draw[cP!80!black, line width=0.8pt] (0,0.04) -- (0.8,0.05) -- (1.1,0.09) -- (1.45,0.17) -- (1.62,0.19) -- (3.6,0.2);
  \node[font=\scriptsize, text=cInk, anchor=west] at (3.62,0.12) {$q(j)$};
\end{scope}
\node[eq] at (0,-0.45) {$A_{t,j} = A_t\, w_{t,j}$,\quad mean$_j\, w_{t,j} = 1$, clipped to $[0.25,4]$\\
  $w_{t,j} \propto q(c_{k+1}) - q(c_k)$ for $j \in [c_k, c_{k+1})$\\
  $q(c) = \pi_\theta(a^\star \mid s_t,\, a_t^{<c},\, \texttt{</think><action>})$};
\end{scope}

% arrows between panels
\draw[-{Stealth[length=2mm]}, cMuted, line width=0.6pt] (5.38,0.3) -- (5.72,0.3);
\draw[-{Stealth[length=2mm]}, cMuted, line width=0.6pt] (11.28,0.3) -- (11.62,0.3);
\end{tikzpicture}
}
\caption{\textbf{Credit at three granularities.} GRPO (left) gives every token of a trajectory one
episode-level advantage. GiGPO and GACA (middle) give each environment step its own advantage, mixing
an anchor-state step contrast $\Ast_t$ with the episode advantage $\Aep$ by a weight $\lambda_t$. All
tokens of the step's response still share it. Probe-weighted credit (right) keeps the step advantage
and redistributes it over the response: the commitment probe $q(j)$ scores how likely the eventual
action is after the first $j$ tokens, and each segment's credit is proportional to how much $q$ rises
across it (weights clipped to $[0.25, 4]$ and renormalized to mean one).}
\label{fig:method}
\end{figure*}
